\section{Cuello de botella de datos: de la anotación manual a la generación de trayectorias impulsada por tutoriales}

\subsection{Por qué los datos de trayectorias son el recurso escaso central}
El presentador señala que las trayectorias de interacción de varios pasos son mucho más costosas que los datos de instrucción de un solo turno. Cada trayectoria de alta calidad suele contener cambios de estado complejos, llamadas a herramientas y recuperación de errores; la recolección puramente manual difícil de sostener el entrenamiento a escala.

\begin{importantbox}{La escasez de datos no está en la ``cantidad de texto'', sino en la ``calidad de la trayectoria ejecutable''}
El valor de los datos de agentes proviene de tres puntos: contexto de tarea completo, acción ejecutable y resultado verificable. Si falta cualquiera de ellos, la ganancia de entrenamiento cae de manera considerable.
\end{importantbox}

\subsection{Construcción de datos impulsada por tutoriales (Tutorial-driven)}
El curso propone extraer plantillas de tareas de alto valor a partir de tutoriales de internet: primero se hace la recolección de tutoriales y el análisis estructurado, luego se convierten en descripciones de tareas estandarizadas y candidatos de pasos, y por último el agente los ejecuta en el entorno para generar las trayectorias.

\begin{figure}[H]
\centering
\includegraphics[width=0.88\textwidth]{frame-05-tutorial-pipeline.jpg}
\caption{Video aproximadamente en 00:44:40: pipeline de construcción de datos impulsada por tutoriales, que incluye extracción, estructuración, ejecución y filtrado}
\end{figure}

\begin{knowledgebox}{Ventajas del esquema impulsado por tutoriales}
\begin{itemize}
    \item Alta diversidad de tareas: cubre la distribución de problemas reales de los usuarios;
    \item Buena transferibilidad: tutoriales del mismo tipo pueden mapearse a tareas de distintos entornos;
    \item Costo controlable: reduce la anotación manual paso a paso;
    \item Fácil de escalar: puede integrarse con scripts de evaluación automática para formar un ciclo cerrado.
\end{itemize}
\end{knowledgebox}

\subsection{Control de calidad: de ``ejecutable'' a ``aprendible''}
El texto de un tutorial no equivale de manera natural a datos de entrenamiento. El curso enfatiza el filtrado en varios niveles: verificación de la ejecutabilidad de la tarea, verificación de la consistencia de la trayectoria, eliminación de pasos erróneos y reducción del peso de muestras con poca información. Los datos con mucho ruido bajan de manera considerable la eficiencia del entrenamiento.

\begin{warningbox}{No confundas ``recolección exitosa'' con ``dato utilizable''}
Muchas trayectorias, aunque se ejecutan por completo, casi no aportan ganancia al aprendizaje del modelo; por ejemplo, rutas de plantilla repetidas, ausencia de pasos de decisión clave o desalineación entre la acción y la observación. La gobernanza de datos debe integrarse al flujo principal de entrenamiento y no ser un parche fuera de línea.
\end{warningbox}

\subsection{Resumen de la sección}
La ruta de datos que propone el curso es ``recolección automática + gobernanza estructurada + verificación por script''. Esta ruta no busca la pureza absoluta, sino aumentar de manera sostenida la densidad de muestras efectivas con un costo controlado.
